\paragraph{Exported fields and their limits.}
Each report holds one frame of nodal values, the last one of the run: coordinates $(r,z)$, logarithmic strains LE11, LE22, LE33, LE12, the von Mises and principal stresses, and the components S11, S22, S33, S12, which in the CAX convention are $\sigma_{rr}$, $\sigma_{zz}$, $\sigma_{\theta\theta}$, $\tau_{rz}$. % src: audit/RAW_AUDIT.md §1; QA_DATA_AND_Z.md §5 (one frame per file)
% FIXED: "final, unloaded frame" → for the completed runs the frame is unloaded (VALIDATION.md §1: median |∫σzz dA|/A = 5.3 MPa on the 1756 cases); for the 158 unfinished runs the drawn part still carries F/A ≈ 232–237 MPa (VALIDATION.md §1). Those runs are removed below.
The targets are read directly from these columns (present in 2443 of 2443 files), and the identity of every stored field is verified against all quantities of the report, not only the four expected ones. % src: QA_DATA_AND_Z.md §2; MANUSCRIPT.md §2; code/verify_pkl.py
The dataset was regenerated from the raw Abaqus reports; the four components $(\sigma_{rr}, \sigma_{\theta\theta}, \sigma_{zz}, \tau_{rz})$ are read from the columns S.S11, S.S33, S.S22 and S.S12 respectively. % src: MANUSCRIPT.md §2 (Data availability); code/raw_probe.py (S.S11 = σ_rr, S.S22 = σ_zz, S.S33 = σ_θθ, S.S12 = τ_rz); code/dataset.py docstring % EDIT: data statement added (critic, E-18)
Nodal averaging leaves a residual at the free surface: on the final dataset the mean $|\sigma_{rr}|$ at $r=R$ is 3.52~MPa and $|\tau_{rz}|$ 0.24~MPa. % src: RESULTS.md header (1756 cases)
This is the accuracy with which the source itself satisfies the traction-free condition and is the baseline for every boundary-condition metric in Section~\ref{sec:res_fem}. % EDIT: \ref{sec:results-physics} → \ref{sec:res_fem} (label defined in sec32)
% REV К1: the enumeration of the quantities absent from the present export, and the reference to a further export, were removed; only the consequence that bears on the present study is kept.
% REV К1 (critic): the framing "one consequence has to be stated explicitly" was dropped — the comment asks for the point to be made less, not more, emphatically; and the statement is confined to the retained cases, so that it does not contradict Section~\ref{sec:res_fem_verification}, where the drawing stress is obtained from the runs whose frame is captured inside the die.
In the retained runs the exported cross-section is unloaded (median $|\int\sigma_{zz}\,\mathrm{d}A|/A$ of 5.3~MPa against an rms $\sigma_{zz}$ of 167~MPa), so the drawing force is not read from these cases. % src: VALIDATION.md §1 (1756 cases: 5.3 MPa against rms 167 MPa)

\paragraph{Axial window.}
Only nodes whose axial coordinate lies within 25--75\,\% of the axial extent of the wire are used (median captured fraction 0.4985); the window is the steady-state part of the drawn wire, of median length 57.5~mm at a radius of 16.2~mm ($L/R \approx 3.5$). % src: audit/RAW_AUDIT.md §2; code/raw_probe.py axial_window; QA_DATA_AND_Z.md §4
% REV К2: the sentence was recast; the shares and the losses are unchanged.
The adequacy of this choice was examined component by component: for each component the share of the in-window variance attributable to the axial direction was measured, together with the loss incurred when the field is collapsed along $z$ (Table~\ref{tab:zshare}). For the normal components these two quantities amount to 10--12\,\% and 2--4\,\% respectively, whereas for $\tau_{rz}$ they reach 40\,\% and 79\,\%: that component changes sign along $z$ and therefore cancels under averaging. % src: QA_DATA_AND_Z.md §4 (shares 0.098–0.116 → 10–12 %; loss 2–4 %; τ_rz 0.397 → 40 %, 79 %)
The shear component is thus intrinsically unpredictable from a one-dimensional representation, a property of the data rather than of any model (Fig.~\ref{fig:fem-fields}: inside the window the mean $|\tau_{rz}|$ of the shown run is 1.06~MPa against 11.0~MPa outside). % src: QA_DATA_AND_Z.md §5 (job 1020_2a_24_rd_1000_cal_50_v_10_fric_050); fig:fem-fields = figures/paper/fig2_fem_fields.png (same job: Q = 0.10, k = 0.5, α = 12°, μ = 0.05, v = 10; caption in manuscript/captions.md Fig. 2) % EDIT: fig:contour-raw → fig:fem-fields, the paper version of the same run (critic); TODO: the figure environment itself is not in any manuscript file yet

% EDIT: coordinate-normalisation note moved from the figure captions into the text (captions.md "Что переезжает в текст"; supervisor's instruction).
% REV К3: the numerical value of the billet radius is not repeated in the text; it is retained in the caption of Fig.~\ref{fig:fem-fields}. % REV К3 (critic): with the value gone, $R$ was left undefined here and clashed with $R$ = surface radius used elsewhere in this subsection; the sentence now says which radius the axis label of Fig.~\ref{fig:fem-fields} denotes.
Coordinates in the field plots are dimensionless throughout. In Fig.~\ref{fig:fem-fields} the radius is normalised by the billet radius, written $R$ in the axis label of that figure, so that the drawn part of the wire ends at $r/R = 0.9$, and the axial coordinate by the length of the model; % src: captions.md Fig. 2 (R = 18 mm; L = 139.5 mm; drawn part ends at r/R = 0.9 at Q = 0.10)
in Figs.~\ref{fig:compare-fields} and \ref{fig:residual-maps} the radius is normalised by the surface radius of the same cross-section, so that $r/R = 1$ is the free surface at every axial position, and $z$ by the window length. % src: captions.md ("Нормировка координат"); code/dataset2d.py (r = RP / RP[:, :, -1])
The radial axis is stretched relative to the axial one: at $L/R \approx 3.5$--$7.7$ an equal-aspect panel degenerates into a strip. % src: captions.md ("Нормировка координат": L/R ≈ 3.5–7.7); experiments/contours.py (set_aspect("auto"))

\begin{table*}[htbp]
\centering\small
\caption{Share of the variance along $z$ inside the axial window and effect of collapsing the field along $z$ by the median (mean absolute value, MPa).} % src: QA_DATA_AND_Z.md §4; TODO: state the population (all 2443 runs or the filtered set) — not stated in the source
\label{tab:zshare}
\begin{tabular}{@{}lccc@{}}
\toprule
Component & Share along $z$ & Before collapse & After collapse \\
\midrule
$\sigma_{rr}$           & 0.114 & 47.9  & 46.6  \\ % src: QA_DATA_AND_Z.md §4
$\sigma_{\theta\theta}$ & 0.116 & 56.6  & 54.4  \\ % src: QA_DATA_AND_Z.md §4
$\sigma_{zz}$           & 0.098 & 140.2 & 137.6 \\ % src: QA_DATA_AND_Z.md §4
$\tau_{rz}$             & 0.397 & 3.22  & 0.68  \\ % src: QA_DATA_AND_Z.md §4
\bottomrule
\end{tabular}
\end{table*}

\paragraph{Profiles and grids.}
For the main track the window nodes are sorted by radius and split into 20 equal-count bins; the bin radius and the four components are the medians within each bin. % src: code/dataset.py (N_R = 20, module docstring); code/raw_probe.py median_bins
Radii are stored per case, and $r$ is normalised by the outermost bin radius of the same case, so that $r=1$ is the free surface by construction (median ratio of the outermost bin to the wire radius 0.996). % src: code/dataset.py (r = r_phys / r_max); audit/RAW_AUDIT.md §2 (0.9958); PLAYBOOK.md §2.3 (0.996)
% REV К5: "nodal noise" replaced by an explicit explanation of the numerical origin of the effect; the figures are unchanged.
Numerical effects inherent in the finite-element treatment of large plastic deformation displace individual nodes in the vicinity of the axis across the axis of symmetry, so that the radial coordinate recovered for them turns out slightly negative ($|r|/R \le 0.054$); in all $1/r$ terms the radius is therefore clipped at $0.05R$. % src: ERRATA.md E-25; code/trainer.py r_clip = 0.05
For the extended track the same nodes are binned on an $8\times20$ grid in $(z,r)$ by the median; $r$ is normalised by the surface radius of the same axial slice, $z$ by the window ($z=0$ entry, $z=1$ exit). % src: code/dataset2d.py (n_z = 8, n_r = 20 defaults; module docstring); QA_DATA_AND_Z.md §4
Targets stay in MPa; since the components differ by two orders of magnitude (mean $|\sigma|$ about 79~MPa for the normal components against 0.67~MPa for $\tau_{rz}$), per-component errors are always reported next to macro-averaged ones. % src: PLAYBOOK.md §2.4

\paragraph{Filtering.}
% REV К6: an explanatory opening sentence was added before the list of the filters.
In the further preparation of the data a part of the computed results had to be discarded, and this was necessary for several distinct reasons. Cases are removed whole, before any split, in the order of Table~\ref{tab:filtering}. The level $Q=0.25$ is excluded entirely: there the process ceased to be drawing in several runs (the billet adhered to the die and the step degenerated into tension), and 17.5\,\% of its runs did not complete against 4.3--6.0\,\% at the other levels. % src: PLAYBOOK.md §5.2 (six "billet adhered to the die" cases + one convergence artefact, all at Q = 0.25); ERRATA.md E-24 (48 of 275 = 17.5 %); TODO: one consistent justification for excluding Q = 0.25 is to be agreed with the corresponding author (code/filters.py labels this filter "elastic regime")
% REV К7: the single run with a mistyped friction label is no longer mentioned in the running text; it is retained as a row of Table~\ref{tab:filtering}. % src: ERRATA.md E-03
Cases whose profile is bit-identical to another case differing only in the friction label are dropped, keeping one per group (323 of the 2443 raw cases, 13.2\,\%; 319 of them survive the two preceding filters and are counted in Table~\ref{tab:filtering}). % src: ERRATA.md E-06; audit/RAW_AUDIT.md §5 (323 / 13.2 % on the raw grid); PLAYBOOK.md §5.5 (−319 after Q = 0.25 and μ-typo removal); code/filters.py find_mu_duplicates
Finally, runs that stopped before the wire had passed through the die are removed: their tail retains the billet radius $R_0$ and the die reference point lies inside the axial extent of the wire (153 of the 159 such runs on the raw grid). % src: ERRATA.md E-24
Where the undrawn part enters the window the surface radius varies along it (spread above 0.1~mm against a median of 0.0026~mm), and what the profile labels as free-surface $\sigma_{rr}$ is contact pressure, down to $-571.3$~MPa against $-19.7$~MPa at worst in clean cases. % src: ERRATA.md E-24 (threshold 0.1 mm → 89 cases; median spread 0.0026 mm; clean 1756 cases min −19.7 MPa)
These 92 cases (89 detected through the window plus 3 with the tail outside it) are removed before splitting: on the random split of the uncleaned set 73 of the 89 fell into the training pool, where the data term contradicts the traction-free condition. % src: ERRATA.md E-24; code/filters.py unfinished_draw_in_window docstring (73 train / 16 test on the 1571/277 random split of the 1848-case set)
% FIXED: "73 of the 89 would otherwise enter training" qualified — the 73/16 count is for the random split of the 1848-case set only.

\begin{table*}[htbp]
\centering\small
\caption{Case-level filtering applied before splitting.}
\label{tab:filtering}
\begin{tabular}{@{}lrr@{}}
\toprule
Step & Removed & Remaining \\
\midrule
Raw reports                                            &     & 2443 \\ % src: audit/RAW_AUDIT.md (header)
Level $Q=0.25$ excluded                                & 275 & 2168 \\ % src: PLAYBOOK.md §5.5; audit/RAW_AUDIT.md §4 (Q = 0.25 → 275)
Friction label outside the design (file-name typo)     & 1   & 2167 \\ % src: PLAYBOOK.md §5.5; ERRATA.md E-03
Bit-identical duplicates differing only in $\mu$ label & 319 & 1848 \\ % src: PLAYBOOK.md §5.5
Runs that did not complete the pass through the die    & 92  & 1756 \\ % src: ERRATA.md E-24; RESULTS_v1_vs_v2.md (1848 → 1756)
\bottomrule
\end{tabular}
\end{table*}

The final dataset has 1756 cases: $35\,120$ radial points in the main track and $280\,960$ grid points in the extended track. % src: 1756×20 and 1756×8×20 (arithmetic); results/runs_v2.csv stage twod: n_train + n_test = 1493 + 263 = 1377 + 379 = 1756, i.e. the same cases
% EDIT: sentence "Removing the last 5.0 % of cases changed the results materially; the comparison with the uncleaned data is given in Section ..." deleted — §3.2 contains no comparison with the uncleaned data and v1 figures are not used (critic)

\paragraph{Group-aware splits.}
% REV К9: an introductory sentence was placed before the description of the partitions.
For the training that follows, several scenarios of partitioning the initial sample were considered, each of them addressing a different question about the behaviour of the surrogates.
% REV К8: the number of points per computation was removed; the uniqueness of the configurations produced by the design and the fact that the partitioning is performed at the level of computations are stated instead.
The runs carried out according to the design of Table~\ref{tab:fem_params} yield, once the filtering described above has been applied, configurations of points that are distinct within the window considered: the computations whose fields coincided were removed as duplicates. The partitioning is accordingly performed at the level of whole computations, so that the points belonging to one computation are never distributed between the training and the test set. % src: code/splits.py sets_to_rows; code/filters.py (module docstring); TODO-group-cv: reference for grouped/blocked hold-out (e.g. Roberts et al. 2017, Ecography 40:913, or scikit-learn GroupKFold)
A random 15\,\% hold-out measures interpolation only; to measure transfer beyond the training range, entire factor levels are held out (Table~\ref{tab:splits}). % src: code/splits.py make_random_split (frac = 0.15); PLAYBOOK.md §7.1
% REV К11: the description of the controls and of the joint-range partition was recast in more formal terms; the content is unchanged.
Each extrapolation split is accompanied by two controls. The \emph{matched} control withholds a random subset of the same size drawn from the whole pool, so that the difference $\Delta\mathrm{RMSE} = \mathrm{RMSE}(\text{extrap}) - \mathrm{RMSE}(\text{matched})$ separates the cost of leaving the training range from the cost of a smaller training set; % src: code/splits.py make_matched_split, module docstring
the \emph{interior} control withholds an inner level of the same axis ($\alpha = 12^\circ$, $Q = 0.10$), which separates the cost of a gap inside the factor grid from the cost of leaving its range. % src: code/splits.py REGIONS (alpha_mid, Q_mid); RESULTS.md §1
% REV К10: the term "corner split" was replaced by a description in terms of a joint departure from the ranges of two factors; $\alpha$ is named the die semi-angle, i.e.\ the half-angle of the deformation zone.
A further partition, denoted \texttt{extrap\_joint:corner} in the run log, is defined by a joint departure from the ranges of two factors, the die semi-angle $\alpha$, i.e.\ the half-angle of the deformation zone, and the reduction $Q$: the region $\alpha = 20^\circ \wedge Q \ge 0.15$ is withheld as a whole. Along either of the two axes taken separately the test cases of this partition remain inside the training range, so that what is probed is a region of the joint factor space left uncovered by the design, and not extrapolation beyond the range of a single factor. % src: code/splits.py REGIONS corner (kind extrap_joint); PLAYBOOK.md §7.1
A coverage report states, for each axis, whether the test range leaves the training range; transfer beyond the training range is claimed only for those axes for which it does. % src: code/splits.py coverage_report
The partitions are fixed (seed 42); the five training seeds of Section~\ref{sec:models} vary the initialisation, the batch order and the 15\,\% early-stopping subset carved from the training pool. % src: code/splits.py defaults (seed = 42); results/runs_v2.csv (seeds 0–4); code/trainer.py (manual_seed, val_frac = 0.15, batch Generator seeded by cfg.seed)
% FIXED: "validation subset" → "early-stopping subset" to keep the word validation for comparison against experiment (VALIDATION.md: verification, not validation).

\begin{table*}[htbp]
\centering\small
\caption{The 13 case-level splits of the 1756 cases. $n_\text{train}$ is the training pool, including the early-stopping subset.} % src: results/runs_v2.csv (n_train, n_test, stage main; 13 distinct splits)
\label{tab:splits}
\begin{tabular}{@{}lllrr@{}}
\toprule
Split & Held-out cases & Kind & $n_\text{train}$ & $n_\text{test}$ \\
\midrule
\texttt{random}              & random 15\,\%                         & interpolation           & 1493 & 263 \\ % src: results/runs_v2.csv
\texttt{extrap:alpha\_max}   & $\alpha = 20^\circ$                   & extrapolation, $\alpha$ & 1377 & 379 \\ % src: results/runs_v2.csv
\texttt{extrap:Q\_high}      & $Q = 0.20$                            & extrapolation, $Q$      & 1484 & 272 \\ % src: results/runs_v2.csv
\texttt{extrap:k\_max}       & $k = 1.0$                             & extrapolation, $k$      & 1466 & 290 \\ % src: results/runs_v2.csv
\texttt{extrap:v\_max}       & $v = 250$~m/min                       & extrapolation, $v$      & 1255 & 501 \\ % src: results/runs_v2.csv
\texttt{extrap\_joint:corner}& $\alpha = 20^\circ \wedge Q \ge 0.15$ & gap in the joint $(\alpha, Q)$ range & 1642 & 114 \\ % REV К10: "joint-space gap" named after the two factors involved % src: results/runs_v2.csv
\texttt{interp:alpha\_mid}   & $\alpha = 12^\circ$                   & interior control        & 1416 & 340 \\ % src: results/runs_v2.csv
\texttt{interp:Q\_mid}       & $Q = 0.10$                            & interior control        & 1438 & 318 \\ % src: results/runs_v2.csv
\texttt{matched:alpha\_max}  & random, same size                     & matched control         & 1377 & 379 \\ % src: results/runs_v2.csv
\texttt{matched:Q\_high}     & random, same size                     & matched control         & 1484 & 272 \\ % src: results/runs_v2.csv
\texttt{matched:k\_max}      & random, same size                     & matched control         & 1466 & 290 \\ % src: results/runs_v2.csv
\texttt{matched:v\_max}      & random, same size                     & matched control         & 1255 & 501 \\ % src: results/runs_v2.csv
\texttt{matched:corner}      & random, same size                     & matched control         & 1642 & 114 \\ % src: results/runs_v2.csv
\bottomrule
\end{tabular}
\end{table*}
